\begin{CVsidebar}

% Added by HSN
\section{Education}
\textbf{University of Tehran}  \\
M.Sc. {\footnotesize in Physics} \\
\softtext{{\footnotesize Tehran, IRAN \\ 
\textit{Sept. 2017 – Sept. 2020}}} \\
\vspace{0.6em}
\textbf{Sharif University of Technology}  \\
B.Sc. {\footnotesize in Physics} \\
\softtext{{\footnotesize Tehran, IRAN \\ 
\textit{Sept. 2011 – Sept. 2016}}}

\section{Skills}

% Hsn added: skills restructured. Three changes:
% Hsn added: (1) \CVrating dropped entirely. "SQL 3/5" and "JavaScript 3/5" from
% Hsn added:     someone who architected a CRM with real-time analytics is a claim made
% Hsn added:     against himself that no reader asked for, and no ATS can read a circle.
% Hsn added:     The macro stays defined in Code/CV.cls for future use. The ~8 lines it
% Hsn added:     freed pay for the new categories.
% Hsn added: (2) A dedicated "AI & LLM Engineering" block, so a recruiter scanning for
% Hsn added:     AI experience finds it in one second and it reinforces the page-2
% Hsn added:     project entries.
% Hsn added: (3) Ordering is the point: Python leads the languages and Django/Laravel/
% Hsn added:     React lead the frameworks, with the WordPress ecosystem kept but demoted
% Hsn added:     below them, so a fast scan reads "engineer" before "WordPress developer".
% Hsn added: Languages Proficiency moved to the bottom of the sidebar so the five skill
% Hsn added: categories read as one uninterrupted block; it used to sit wedged between
% Hsn added: Programming Languages and Backend Development.
\textbf{Programming Languages:} \\
\vspace{0.1em}
{\small Python, PHP, JavaScript/TypeScript, C\#, SQL} \\
\dotline

\textbf{AI \& LLM Engineering:} \\
\vspace{0.1em}
{\small OpenAI API, LLM integration, prompt engineering, chat-interface design, AI-assisted development (Claude Code, GitHub Copilot)} \\
\dotline

\textbf{Backend Development:} \\
\vspace{0.1em}
{\small Django, Laravel, Flask, PostgreSQL, MySQL, Redis, MongoDB, REST API/GraphQL, JWT/OAuth2, WordPress, WooCommerce, Composer, Domain-Driven Design (DDD), Microservices architecture} \\
\dotline

\textbf{Frontend Development:} \\
\vspace{0.1em}
{\small React, Next.js, Tailwind, Livewire, Inertia.js, Blade, Alpine.js, jQuery, Bootstrap, Elementor, Accessibility (a11y)} \\
\dotline

\textbf{Systems \& Workflow:} \\
\vspace{0.1em}
{\small Git, Docker, Linux, Apache, CI/CD, Caching strategies, SOLID, Role-Based Access Control (RBAC), High Performance Computing (HPC)} \\
\dotline

% Hsn added: Arabic (Elementary) cut -- "Elementary" tells a reader nothing actionable,
% Hsn added: while a documented TOEFL score is real de-risking evidence for the remote
% Hsn added: and relocation applications. Also fixed the casing of TOEFL.
\textbf{Languages Proficiency:} \\
\vspace{0.1em}
English \hfill \softtext{\textit{Fluent}{\footnotesize \textit{(TOEFL: 102)}}} \\
Farsi \hfill \softtext{\textit{Native}} \\

%  Kubernetes, Nginx, Apache, AWS, DigitalOcean, Azure, Database indexing and optimization,
%  Caching strategies, JWT / OAuth2, Microservices, Domain-Driven Design (DDD), SOLID, TDD / BDD, Agile / Scrum, Documentation & API design,
%   Code review, System design, Security best practices (CSRF, XSS, SQL injection prevention), Role-Based Access Control (RBAC), 
%   Encryption and hashing (bcrypt, Argon2)


% \textbf{DevOps & Tools:} \\
% \vspace{0.1em}
% Git, Docker, Linux, MS Windows \\
% \dotline








% \section{About me}
% Age: \softtext{XX years old, born DD.MM.YYYY} \par
% Interrests: \softtext{Something}

% \vspace{\baselineskip}
% \begin{quote}
%     My favourite quote.
% \end{quote}
% NOTE: you can change the quote colour 
% by adding a colour in square brackets [ ]:
% \begin{quote}[red] will make the quote
% appear in red. Similarly for soft_text and
% other colours defined in the preamble or .cls

% \section{Languages}

% % Options for displaying languages:
% %-------------------------------
% % * \LngLvlCEFR{<Language>}{<Level>}:
% %      Specify CEFR language <Level>: 
% %      A1, A2, B1, B2, C1, C2 or 0
% %
% % * \LngLvlILR{<Language>}{<Level>}:
% %       Specify ILR language <Level>:
% %       0, 1, 2, 3, 4 or 5
% %
% % * \CVrating[<text>]{<Language>}{<level>}{<max level>}:
% %       <text> will appear before rating circles
% %       it will display <level>/<max level> circles,
% %       i.e. if you specify 2 and 3 it will show 
% %       2 green circles out of a total of 3 circles.
% %
% % * \LngLvlWords{<Language>}{<text>}:
% %       appears as <Language> \hfill \textit{<text>}
% %       good for expressing language proficiency with words
% %       i.e. " English      Professional Working knowledge "
% %       or   " English                   Intermidiate User "
% %
% % * \LngTextCFR :
% %       prints text:
% %       " *Using CEFR rating for language proficency "
% %       in current document language
% %
% % * \LngTextILR :
% %       prints text:
% %       " *Using ILR rating for language proficency "
% %       in current document language
% %
% % Comment: Preferably only use one style or you may combine
% %       \LngLvlCEFR / \LngLvlILR with \LngLvlWords in  this manner:
% %       \LngLvlCEFR{<language>}{<Level>} and \LngLvlWords{}{<description>}

% \LngLvlCEFR{Norwegian}{C2}
% \LngLvlWords{}{Native}
% \LngLvlCEFR{Spanish}{A2}
% \LngTextCEFR

% \LngLvlILR{English}{4}
% \LngTextILR

% \LngLvlWords{Language}{Basic Knowledge}

% %% For unsupported languages for \LngTextCEFR and \LngTextILR, use this:
% % \vspace{0.5\baselineskip}
% % {\footnotesize\softtext{*Using CEFR rating for language proficency.}}

% \dotline

% \CVrating{Python}{3}{5}
% \CVrating[Pro]{\LaTeX}{6}{6}
% C++

% \section{Certificates}
% % You can use itemlist instead of itemize.
% % This has less space between items.
% \begin{itemlist}
%     \item Driver's licence; B, BE, T
%     \item Chainsaw, 2017
% \end{itemlist}

% \section{Referees}

% % e-mails will get hyperref. 
% \CVref{Person 1}{example@email.com}{Adress}{Institute}

% \CVref{Person 2}{example@email.com}{Adress}{Institute}



\end{CVsidebar}